\documentclass[12pt]{article}
\usepackage[margin=0.75in,top=0.8in,bottom=0.65in]{geometry}
\usepackage{lmodern}
\usepackage{microtype}
\usepackage{fancyhdr}
\usepackage{titlesec}
\usepackage{enumitem}
\usepackage{lastpage}
\usepackage[hidelinks]{hyperref}

\setlength{\parindent}{0pt}
\setlength{\parskip}{0.38em}
\setlength{\headheight}{26pt}
\titleformat{\section}{\large\bfseries\scshape}{}{0pt}{}
\titleformat{\subsection}{\normalsize\bfseries}{}{0pt}{}
\pagestyle{fancy}
\fancyhf{}
\fancyhead[C]{\footnotesize Lit Example Reader \quad Assignment ID: U1L3LE \quad Page \thepage\ of \pageref{LastPage}}
\renewcommand{\headrulewidth}{0.5pt}

\newenvironment{playpassage}{\begin{quote}\small\setlength{\parskip}{0.25em}}{\end{quote}}
\newcommand{\speaker}[1]{\textbf{#1.}}

\begin{document}

\begin{center}
{\LARGE\bfseries Week 3 Lit Example Reader}\par\vspace{0.05em}
{\large\scshape Macbeth: Prophecies as Propositions}\par\vspace{0.15em}
\rule{0.78\textwidth}{0.5pt}
\end{center}

\textbf{Purpose:} This reader gives Macbeth passages for applying Week 3's logic habit: distinguish propositions from other utterances, translate ordinary language into clear categorical form, and notice quantity, quality, and distribution.

\textbf{What to watch for:} Macbeth treats some spoken lines as if they were fixed logical guarantees. But not every powerful line is a clear proposition, and not every proposition has the same quantity. Before reasoning from a line, ask exactly what kind of claim it is making.

\section*{Before the Passages: The Week 3 Logic Tool}

A proposition is a claim that can be treated as true or false. In categorical logic, a clear proposition has a subject, a predicate, a quality, and a quantity:

\begin{itemize}[leftmargin=1.5em,itemsep=0.2em]
\item \textbf{Subject:} what the claim is about.
\item \textbf{Predicate:} what is being affirmed or denied of the subject.
\item \textbf{Quality:} affirmative or negative.
\item \textbf{Quantity:} universal or particular.
\end{itemize}

The four standard forms are A (All S are P), E (No S are P), I (Some S are P), and O (Some S are not P). Translating a dramatic line into standard form is not the same as proving it true. Translation only clarifies what the claim says.

\section*{Passage 1: Act 4, Scene 1 --- The Apparitions Speak}

\textbf{Context:} Macbeth returns to the witches and demands certainty. The apparitions give him commands and claims. Logic begins by sorting which lines are propositions and what their form is.

\begin{playpassage}
\speaker{First Apparition} Macbeth! Macbeth! Macbeth! beware Macduff;\par
Beware the thane of Fife. Dismiss me. Enough.

\speaker{Second Apparition} Be bloody, bold, and resolute; laugh to scorn\par
The power of man, for none of woman born\par
Shall harm Macbeth.

\speaker{Third Apparition} Be lion-mettled, proud; and take no care\par
Who chafes, who frets, or where conspirers are:\par
Macbeth shall never vanquish'd be until\par
Great Birnam wood to high Dunsinane hill\par
Shall come against him.
\end{playpassage}

\subsection*{Reader Notes}

\textit{Beware Macduff} is a command or warning, not a categorical proposition by itself. \textit{None of woman born shall harm Macbeth} is closer to a universal negative claim: no persons born of woman are persons who will harm Macbeth. The Birnam Wood line is a conditional or time-limited prediction, not a simple A/E/I/O proposition unless it is translated carefully.

\textbf{Reading task:} Mark each line as command, warning, categorical claim, or conditional prediction.

\newpage

\section*{Passage 2: Act 5, Scenes 5 and 8 --- Macbeth Tests the Claims}

\textbf{Context:} Macbeth later learns that Birnam Wood appears to be moving because soldiers carry branches. He still leans on the other proposition: no one born of woman can harm him. Then Macduff clarifies a fact that changes the subject class.

\begin{playpassage}
\speaker{Messenger} As I did stand my watch upon the hill,\par
I look'd toward Birnam, and anon, methought,\par
The wood began to move.

\speaker{Macbeth} I pull in resolution, and begin\par
To doubt the equivocation of the fiend\par
That lies like truth: ``Fear not, till Birnam wood\par
Do come to Dunsinane:'' and now a wood\par
Comes toward Dunsinane.

\speaker{Macbeth} Thou losest labour:\par
As easy mayst thou the intrenchant air\par
With thy keen sword impress as make me bleed:\par
Let fall thy blade on vulnerable crests;\par
I bear a charmed life, which must not yield,\par
To one of woman born.

\speaker{Macduff} Despair thy charm;\par
And let the angel whom thou still hast served\par
Tell thee, Macduff was from his mother's womb\par
Untimely ripp'd.
\end{playpassage}

\subsection*{Reader Notes}

Macbeth's confidence depends on a proposition about a class: those \textit{born of woman}. If the subject class is misunderstood, the inference becomes dangerous. Macduff does not merely boast. He supplies a fact that forces Macbeth to ask whether Macduff belongs inside or outside the class named by the proposition.

\textbf{Reading task:} Underline the phrase \textit{of woman born}. Circle Macduff's line that challenges whether he belongs in that class.

\section*{How to Use These Passages on the Worksheet}

The worksheet will ask you to translate dramatic language into logical form. Do not turn every line into the same kind of claim. Commands, predictions, categorical propositions, and conditional statements behave differently.

\textbf{Strong college-level answers} should do three things:

\begin{itemize}[leftmargin=1.5em,itemsep=0.2em]
\item say whether a line is a proposition before classifying it;
\item identify subject, predicate, quality, and quantity when a categorical form is possible;
\item explain how translation affects Macbeth's reasoning without pretending that translation proves the prophecy true.
\end{itemize}

\end{document}
